% bracket_ink_issue105.tex - 回归测试: 直排夹注号的字面分布
% Issue: https://github.com/open-guji/luatex-cn/issues/105 （第 ② 条）
%
% 原始报告：「结束夹注符号的位置偏上」——》 骑到前一个字上。
%
% 直排下引擎按字形的 height/depth 盒把字面居中在字幅里，而 vert 形的
% height/depth 常常量不准（墨迹整个在基线之上时 depth 被截为 0），字面
% 因此随字体上下漂：同一段文字在思源宋体里 》 贴前字、在京華老宋体里
% 却是居中。现按 clreq 6.3.2.1 图 30 锚定：
%   开始夹注符号（︽︵﹁）字面在下半格，贴后文；
%   结束夹注符号（︾︶﹂）字面在上半格，贴前文。
%
% 图 30 的「字面上下两侧（居中）」一组只有点号与间隔号，夹注号在中国大陆式
% 与台湾式下都是单侧——两种风格的夹注号字面应当**落在同一处**。本用例因此
% 同页两栏同文：上栏台湾式、下栏中国大陆式，只切 punct-style，版式不变，
% 单看这一个轴。对照物是点号：它才是分风格的（台湾式居中、中国大陆式偏靠
% 右上），两栏的句号必然不同位。
%
% 注：一页内每个「正文」块自带一套标点风格，但块与块的列是各自从头排的，
% 所以两栏靠 分栏=2 + 换栏 错开，不能直接写两个「正文」（会重叠）。
\documentclass{ltc-tw-vbook}
\setmainfont{TW-Kai}
\begin{document}

% ---- 上栏：台湾式（本文档类默认） ----
\begin{正文}[分栏=2, 栏间距=16pt]

台湾式：夹注号单侧、点号居中。

书名号：推荐《史记》这本书。篇名号：〈项羽本纪〉。

括号：排版（版面设计）是手艺。方括号：【要点】需留意。

引号：他说「这本书很好」。嵌套：「我读过『论语』。」

连续夹注号：（甲）「乙」《丙》三处相邻。

\end{正文}

% ---- 下栏：中国大陆式（只切 punct-style，其余不动） ----
\标点设置{punct-style=mainland}

\begin{正文}[分栏=2, 栏间距=16pt]
\换栏

中国大陆式：夹注号同上，点号改为偏靠右上。

书名号：推荐《史记》这本书。篇名号：〈项羽本纪〉。

括号：排版（版面设计）是手艺。方括号：【要点】需留意。

引号：他说「这本书很好」。嵌套：「我读过『论语』。」

连续夹注号：（甲）「乙」《丙》三处相邻。

\end{正文}
\end{document}
